\documentclass[11pt]{article}
% ===========================================================================
%  Shared preamble for the AP Calculus AB exam-preparation worksheets.
%  Input right after \documentclass. Files beginning with "_" are not compiled
%  on their own. Navy and gold, to match the course's lesson pages.
% ===========================================================================
\usepackage[letterpaper,top=2.3cm,bottom=2.2cm,left=1.7cm,right=1.7cm,headheight=28pt,footskip=30pt]{geometry}
\usepackage{amsmath,amssymb}
\usepackage[T1]{fontenc}
\usepackage[utf8]{inputenc}
\usepackage{lmodern}
\usepackage{xcolor}
\usepackage[most]{tcolorbox}
\usepackage{enumitem}
\usepackage{fancyhdr}
\usepackage{array}
\usepackage{tikz}
\usepackage{pgfplots}
\pgfplotsset{compat=1.18}

\definecolor{navy}{HTML}{161E54}
\definecolor{ink}{HTML}{27357E}
\definecolor{gold}{HTML}{B98A00}
\definecolor{goldbg}{HTML}{FFFDF3}
\definecolor{inkbg}{HTML}{F6F7FD}
\definecolor{rule}{HTML}{C9CFF0}
\definecolor{okgreen}{HTML}{1B7A44}
\definecolor{okbg}{HTML}{F1F8F3}

\pagestyle{fancy}
\fancyhf{}
\fancyhead[L]{\footnotesize Name: \underline{\hspace{4cm}}}
\fancyhead[C]{\textbf{\color{navy}Integration Academy}\\[-2pt]{\scriptsize Calculus AB \textperiodcentered{} AP\textsuperscript{\textregistered} Exam Preparation}}
\fancyfoot[L]{\tiny\color{black!55}\parbox[t]{0.78\textwidth}{AP\textsuperscript{\textregistered} is a registered trademark of the College Board, which was not involved in the production of, and does not endorse, this material. \textcopyright{} Integration Academy.}}
\fancyfoot[R]{\scriptsize\thepage}
\renewcommand{\headrulewidth}{1.2pt}
\renewcommand{\headrule}{\hbox to\headwidth{\color{gold}\leaders\hrule height \headrulewidth\hfill}}
\renewcommand{\footrulewidth}{0pt}
\setlength{\parindent}{0pt}
\setlength{\parskip}{4pt}

% Title block: \wstitle{code}{title}{unit line}
\newcommand{\wstitle}[3]{\begin{center}{\small\bfseries\color{gold}WORKSHEET #1}\\[3pt]{\LARGE\bfseries\color{navy}#2}\\[3pt]{\small\color{black!60}#3}\end{center}\vspace{2pt}}
% Section heading
\newcommand{\sech}[1]{\par\vspace{10pt}{\large\bfseries\color{navy}#1}\par\nobreak\vspace{2pt}{\color{rule}\hrule height0.8pt}\vspace{7pt}}
% A level tag printed beside a group of questions
\newcommand{\level}[1]{\par\vspace{6pt}{\footnotesize\bfseries\color{gold}\MakeUppercase{#1}}\par\nobreak\vspace{2pt}}

\newtcolorbox{ideabox}{colframe=ink,colback=inkbg,coltitle=white,colbacktitle=ink,fonttitle=\bfseries\small,title={Key ideas},arc=2pt,boxrule=0.8pt,left=8pt,right=8pt}
\newtcolorbox{example}[1]{colframe=rule,colback=white,coltitle=navy,colbacktitle=inkbg,fonttitle=\bfseries\small,title={#1},arc=2pt,boxrule=0.8pt,borderline west={3pt}{0pt}{ink},left=10pt,right=8pt}
\newtcolorbox{frbox}[1]{colframe=gold,colback=goldbg,coltitle=white,colbacktitle=gold,fonttitle=\bfseries\small,title={#1},arc=2pt,boxrule=0.8pt,breakable,left=8pt,right=8pt}
\newtcolorbox{solbox}{colframe=okgreen,colback=okbg,coltitle=white,colbacktitle=okgreen,fonttitle=\bfseries\small,title={Worked solutions},arc=2pt,boxrule=0.8pt,breakable,left=8pt,right=8pt}
\newtcolorbox{scorebox}[1]{colframe=navy,colback=white,coltitle=white,colbacktitle=navy,fonttitle=\bfseries\small,title={#1},arc=2pt,boxrule=0.8pt,breakable,left=8pt,right=8pt}

\newcommand{\soln}{\par\smallskip\textbf{\color{ink}Solution.}\ }
\newcommand{\result}[1]{\par\smallskip\textbf{#1}}
% blank working space of a given height
\newcommand{\work}[1]{\par\vspace{#1}}
% a point value in the margin of a scoring guide
\newcommand{\pts}[1]{\hfill{\footnotesize\bfseries\color{gold}[#1]}}

% Coordinate grid for graph questions: \begin{apgraph}{xmin}{xmax}{ymin}{ymax} ... \end{apgraph}
\newenvironment{apgraph}[4]{\begin{center}\begin{tikzpicture}\begin{axis}[
  axis lines=middle, xmin=#1, xmax=#2, ymin=#3, ymax=#4,
  grid=both, grid style={black!12}, xtick distance=1, ytick distance=1,
  ticklabel style={font=\scriptsize}, xlabel={$x$}, ylabel={$y$},
  width=10cm, height=7.2cm, samples=80, clip=false,
  every axis plot/.append style={line width=1.1pt, color=ink}]}{\end{axis}\end{tikzpicture}\end{center}}
\newcommand{\closeddot}[2]{\addplot[only marks, mark=*, mark size=2.2pt, color=ink] coordinates {(#1,#2)};}
\newcommand{\opendot}[2]{\addplot[only marks, mark=*, mark size=2.2pt, color=ink, mark options={fill=white}] coordinates {(#1,#2)};}

% Multiple-choice options: four across, or stacked when the options are long.
\newcommand{\opts}[4]{\par\vspace{3pt}\begin{tabular}{@{}p{0.235\linewidth}p{0.235\linewidth}p{0.235\linewidth}p{0.235\linewidth}@{}}(A)\ #1 & (B)\ #2 & (C)\ #3 & (D)\ #4\end{tabular}\par\vspace{9pt}}
\newcommand{\optsv}[4]{\par\vspace{3pt}\begin{tabular}{@{}l@{\ \ }p{0.9\linewidth}@{}}(A) & #1\\(B) & #2\\(C) & #3\\(D) & #4\end{tabular}\par\vspace{9pt}}

\fancyhead[R]{\footnotesize\color{black!60}Worksheet 2.4}
\begin{document}

\wstitle{2.4}{Corners, Cusps and Breaks: When the Derivative Fails}{Unit 2: The Derivative}

\begin{ideabox}
\begin{itemize}[leftmargin=*,itemsep=2pt,topsep=2pt]
  \item A function is not differentiable at a \textbf{corner}, a \textbf{cusp}, a \textbf{vertical tangent} or a \textbf{discontinuity}.
  \item Differentiable implies continuous. Continuous does \emph{not} imply differentiable.
  \item A piecewise function is differentiable at a joining point when the pieces \textbf{meet} (continuity) and have the \textbf{same slope} there.
  \item To find unknown constants, write one equation for continuity and one for equal slopes.
\end{itemize}
\end{ideabox}

\sech{Worked examples}
\begin{example}{Example 1 \textperiodcentered{} A corner}
Show that $f(x)=|x+2|$ is not differentiable at $x=-2$.\soln For $h>0$, $\dfrac{|h|}{h}=1$; for $h<0$, $\dfrac{|h|}{h}=-1$. The one-sided limits of the difference quotient are $1$ and $-1$.\result{The limit does not exist, so $f'(-2)$ does not exist.}
\end{example}

\begin{example}{Example 2 \textperiodcentered{} Checking a join}
Is $f(x)=\begin{cases}x^2,&x\le3\\6x-9,&x>3\end{cases}$ differentiable at $x=3$?\soln Values: $9$ and $18-9=9$, so $f$ is continuous. Slopes: $2(3)=6$ and $6$.\result{Yes: the pieces meet with the same slope.}
\end{example}

\begin{example}{Example 3 \textperiodcentered{} Finding constants}
Find $a$ and $b$ so that $f(x)=\begin{cases}ax+b,&x<2\\x^2+1,&x\ge2\end{cases}$ is differentiable at $x=2$.\soln Slopes: $a=2(2)=4$. Continuity: $2a+b=5$.\result{$a=4$ and $b=-3$.}
\end{example}

\newpage
\sech{Practice}

\level{Core}
\begin{enumerate}[leftmargin=*,itemsep=2pt]
  \item At what value of $x$ is $f(x)=|x-5|$ not differentiable?\work{2.3cm}
  \item Is $f(x)=x^{2/3}$ differentiable at $x=0$? Give a reason.\work{2.3cm}
  \item Is $f(x)=\begin{cases}4x-1,&x\le1\\2x^2+1,&x>1\end{cases}$ differentiable at $x=1$?\work{2.3cm}
  \item True or false, with a reason: every continuous function is differentiable.\work{2.3cm}
\end{enumerate}

\level{Exam level}
\begin{enumerate}[leftmargin=*,itemsep=2pt,start=5]
  \item Is $f(x)=\begin{cases}x^2+2x,&x\le0\\\sin(2x),&x>0\end{cases}$ differentiable at $x=0$? Justify your answer.\work{2.8cm}
  \item Find $a$ and $b$ so that $f(x)=\begin{cases}ax^2,&x\le2\\4x+b,&x>2\end{cases}$ is differentiable at $x=2$.\work{2.8cm}
  \item At what values of $x$ is $f(x)=\left|x^2-1\right|$ not differentiable?\work{2.8cm}
  \item Is $f(x)=\begin{cases}\sqrt{x},&0\le x\le4\\\dfrac x4+1,&x>4\end{cases}$ differentiable at $x=4$? Justify your answer.\work{2.8cm}
\end{enumerate}

\begin{frbox}{Question 9 \textperiodcentered{} Free-response practice}
Let $g(x)=\begin{cases}x^2-4x,&x\le1\\ax+b,&x>1\end{cases}$ where $a$ and $b$ are constants.
\begin{enumerate}[label=(\alph*),leftmargin=*,itemsep=2pt]
  \item Write an equation in $a$ and $b$ that makes $g$ continuous at $x=1$.\work{2.6cm}
  \item Find the values of $a$ and $b$ for which $g$ is differentiable at $x=1$.\work{2.6cm}
  \item With those values, write an equation of the line tangent to the graph of $g$ at $x=1$. What do you notice?\work{2.6cm}
\end{enumerate}
\end{frbox}

\newpage
\begin{solbox}
\begin{enumerate}[leftmargin=*,itemsep=4pt]
  \item $x=5$.
  \item No. $f'(x)=\dfrac23x^{-1/3}$ is unbounded near 0: the graph has a cusp.
  \item Yes. Values: $3$ and $3$. Slopes: $4$ and $4(1)=4$.
  \item False. $f(x)=|x|$ is continuous everywhere but not differentiable at 0.
  \item Yes. Values: $0$ and $\sin0=0$. Slopes: $2(0)+2=2$ and $2\cos0=2$.
  \item Slopes: $4a=4$, so $a=1$. Continuity: $4=8+b$, so $b=-4$.
  \item $x=-1$ and $x=1$, where $x^2-1$ changes sign and the graph has corners.
  \item Yes. Values: $2$ and $2$. Slopes: $\dfrac{1}{2\sqrt4}=\dfrac14$ and $\dfrac14$.
  \item \begin{enumerate}[label=(\alph*),leftmargin=*,itemsep=3pt]
    \item $1-4=a+b$, so $a+b=-3$.
    \item Slopes: $2(1)-4=-2=a$. Then $b=-1$.
    \item The point is $(1,-3)$ and the slope is $-2$: $y+3=-2(x-1)$, that is $y=-2x-1$. This is the linear piece itself: the right-hand piece is the tangent line to the parabola.
  \end{enumerate}
\end{enumerate}
\end{solbox}

\end{document}
